Transformer architectures for long-horizon forecasting \cite{zhou2020informer,wu2021autoformer} reported large gains on the ETT and related benchmarks. Zeng et al.\ \cite{zeng2022are} then showed that a one-layer linear model with a moving-average decomposition (DLinear) outperformed these models on most of the same benchmarks, which cast doubt on the benefit of temporal attention. Patching and channel independence \cite{nie2022time} restored a Transformer advantage in later reports, while benchmark studies found that evaluation details matter and that simple methods are often competitive \cite{qiu2024tfb}. Each real dataset mixes trend, seasonalities, noise and non-stationarity in unknown proportions, so the question ``are simple linear models enough?'' has no single answer on real data.

This paper asks a narrower, controlled question: \emph{on series whose trend, seasonal structure, noise level and regime switching we set ourselves, which property, if any, lets a small nonlinear model beat a DLinear-style model at horizons 24, 96 and 192?} We reimplement four systems (seasonal naive, DLinear-style, PatchTST-style, patch GRU), give them the same look-back, normalisation, optimiser and budget, and test five pre-registered hypotheses (Section~\ref{sec:setup}) with an explicit win rule. We also run ETTh1 \cite{zhou2020informer} as one real series, two sweeps (regime length and noise level), a design ablation and a training-budget sweep.

Contributions: (i) a protocol-matched comparison across six tasks and three horizons with three seeds each; (ii) a budget-dependent answer: at 400 training steps no nonlinear model passes the win rule anywhere in the main grid, whereas at 1000 and 2000 steps the nonlinear models pass it on the regime-switching task at horizons 24 and 96 (the GRU at both, the patch transformer at 24) and on no other task that we tested; (iii) sweeps showing that this margin is absent at horizon 192 and for long regimes, and that the low-budget patch-transformer margin depends on noise; (iv) code, per-run logs and the run registry. The budget sweep was not in the original plan; it was added after an independent audit showed that the 400-step verdicts did not hold at 2000 steps. All claims are limited to the small scale described in Section~\ref{sec:limitations}.
